\section*{Informatik}

\subsection*{Willkommen!}
Willkommen in der vielfältigen Welt der Informatik! \emoji{robot} \emoji{smile}

Ich freue mich sehr darauf, Ihnen die vielen Aspekte dieser spannenden Wissenschaft aufzuzeigen. In den kommenden drei Semestern werden Sie die vielfältigen Facetten der digitalen Welt und der Computer besser kennenlernen. 

Im ersten Block lernen Sie die \textbf{Grundlagen der Informatik} kennen:
\begin{itemize}
\item \textbf{Programmieren}: Wie können wir dem Computer spannende und komplexe Aufgaben geben, welche unseren Alltag erleichtern und unsere Kreativität beflügeln?
\item \textbf{Netzwerke}: Wie funktioniert das Internet?
\item \textbf{Zahlensysteme}: Wie funktionieren Computer? 
\end{itemize}

In einem zweiten Block sprechen wir über Data Science und Algorithmen, also Programm-Abläufe, die unser Leben massgeblich prägen:
\begin{itemize}
\item \textbf{Verschlüsselung}: Wie können wir unsere Daten schützen, so dass keine unbefugten Personen darauf Zugriff haben?
\item \textbf{Kompression}: Wie können wir die riesigen Datenberge, die jeden Tag wachsen, effizient speichern?
\item \textbf{Selbstkorrigierende Codes}: Wie können wir sicherstellen, dass es keine Fehler bei der Übertragung von Daten gibt?
\item \textbf{Datenbanken}: Wie können wir effizient auf Daten zugreifen, diese verändern und zusammenfügen, um spannende neue Einsichten zu gewinnen?
\item \textbf{Machine Learning und Klassifikation}: Wie können wir neue Einsichten aus grossen Datenmengen gewinnen? Wie können wir dem Computer beibringen, automatische Vorhersagen für neue Daten zu treffen (Wetter, Pflanzenart, etc.)?
\end{itemize}

Nicht zuletzt werden wir während des gesamten Unterrichts immer wieder über gesellschaftliche Folgen der Digitalisierung sprechen und neuste Trends kritisch gemeinsam betrachten. Dieses Thema werden wir im letzten Block gemeinsam vertiefen.


\subsection*{Grundwerte}
Ich wünsche mir eine \textbf{entspannte} und \textbf{freundliche} Klassen-Atmosphäre, in der sich alle wohl fühlen. Dazu gehört für mich eine Prise (Selbst-)Ironie und \textbf{Humor}, allerdings auch \textbf{Ernsthaftigkeit}, \textbf{Neugier} und \textbf{Offenheit}, um sich auf Neues einzulassen. 

Wichtig sind mir zudem \textbf{Pünktlichkeit und Verlässlichkeit}: Ich wünsche mir, dass die (wenigen) erteilten Hausaufgaben pflichtbewusst erledigt und Termine notiert werden. Auch ich werde mich an diesem Massstab messen und werde versuchen, Ihre Prüfungen stets zeitnah zu korrigieren und bewerten.

Ich schätze Ihr \textbf{Feedback} auf meinen Unterricht sehr, da mir dies hilft, mich stetig zu verbessern. Daher werde ich Sie regelmässig um ihr freiwilliges, anonymes Feedback zum Unterricht bitten. Sie können mir stets auch persönliche Rückmeldungen geben. Ich freue mich über jegliches konstruktives und möglichst konkret umsetzbares Feedback von Ihnen!

Falls Sie weitere Themen haben, die sie interessieren, dürfen Sie sich sehr gerne einbringen!

Allgemein setze ich alles daran, den Unterricht für Sie möglichst spannend und kurzweilig zu machen (ich kann jedoch nicht garantieren, dass dies immer gelingt).

\subsection*{Regeln}

\begin{itemize}
	\item Bitte bringen Sie in jede Unterrichtsstunde ihren geladenen Laptop mit, sowie ein Ladekabel dafür. 
	\item Ich bitte Sie freundlich, während dem Unterricht nur Wasser zu trinken und Essen sowie andere Getränke auf die Pause zu verschieben
	\item WC sowie Wasser auffüllen nur in begründeten Ausnahmefällen.
	\item Ich bitte Sie, jeweils pünktlich im Unterricht zu erscheinen.
\end{itemize}